% Team 12, Paper 15: Lou & Shu (2017). Slide text is written by hand; every number is \V{key}, a value
% that src/build_docs.py writes into gen/values.tex from the CSVs under results/ (the same keys as in
% README.md and PLAN.md: `python -I src/build_docs.py --list`). Tables: gen/slide8.tex, gen/slide9.tex,
% gen/slide10.tex, gen/sims.tex. Build: `make slides` (xelatex). Plan and speaker notes: PLAN.md.
\documentclass[aspectratio=169,10pt]{beamer}
\usepackage{fontspec}
\setsansfont{Arial}
\usepackage{booktabs}
\usepackage{array}
\usepackage{graphicx}
\usepackage{hyperref}

% generated by src/build_docs.py from results/: do not edit

\expandafter\def\csname v@an.V1m\endcsname{−1.64 (t −0.44)}
\expandafter\def\csname v@an.V1m.var\endcsname{DIF}
\expandafter\def\csname v@an.V1y\endcsname{14.11 (t 1.55)}
\expandafter\def\csname v@an.V1y.var\endcsname{DIF}
\expandafter\def\csname v@an.V2m\endcsname{5.18 (t 1.70)}
\expandafter\def\csname v@an.V2m.var\endcsname{DIF}
\expandafter\def\csname v@an.V2y\endcsname{10.78 (t 1.51)}
\expandafter\def\csname v@an.V2y.var\endcsname{DIF}
\expandafter\def\csname v@an.V3m\endcsname{−6.33 (t −1.98)}
\expandafter\def\csname v@an.V3m.var\endcsname{ResLnA}
\expandafter\def\csname v@an.V3y\endcsname{−2.33 (t −0.39)}
\expandafter\def\csname v@an.V3y.var\endcsname{RlnILLIQ}
\expandafter\def\csname v@an.V4y\endcsname{45.12 (t 2.14)}
\expandafter\def\csname v@an.V4y.var\endcsname{fDIF\_int}
\expandafter\def\csname v@an.V5m\endcsname{−0.143 (t −0.05)}
\expandafter\def\csname v@an.V5m.var\endcsname{ResLnA}
\expandafter\def\csname v@an.V5y\endcsname{3.00 (t 0.56)}
\expandafter\def\csname v@an.V5y.var\endcsname{RlnILLIQ}
\expandafter\def\csname v@an.V6m\endcsname{0.814 (t 1.46)}
\expandafter\def\csname v@an.V6m.var\endcsname{lnA\_C}
\expandafter\def\csname v@an.V6y\endcsname{0.909 (t 1.01)}
\expandafter\def\csname v@an.V6y.var\endcsname{lnIDVOL}
\expandafter\def\csname v@an.c1\endcsname{−8.10 (t −2.37)}
\expandafter\def\csname v@an.c4\endcsname{−0.143 (t −0.05)}
\expandafter\def\csname v@an.cnt.ret.AN_1.96\endcsname{2}
\expandafter\def\csname v@an.cnt.ret.AN_3\endcsname{0}
\expandafter\def\csname v@an.cnt.ret.LS_1.96\endcsname{1}
\expandafter\def\csname v@an.cnt.ret.LS_3\endcsname{0}
\expandafter\def\csname v@an.cnt.ret.LS_null_1.96\endcsname{6}
\expandafter\def\csname v@an.cnt.ret.none_1.96\endcsname{0}
\expandafter\def\csname v@an.cnt.ret_adj.AN_1.96\endcsname{1}
\expandafter\def\csname v@an.cnt.ret_adj.AN_3\endcsname{0}
\expandafter\def\csname v@an.cnt.ret_adj.LS_1.96\endcsname{1}
\expandafter\def\csname v@an.cnt.ret_adj.LS_3\endcsname{0}
\expandafter\def\csname v@an.cnt.ret_adj.LS_null_1.96\endcsname{7}
\expandafter\def\csname v@an.cnt.ret_adj.none_1.96\endcsname{0}
\expandafter\def\csname v@an.corr_res_absret\endcsname{0.78}
\expandafter\def\csname v@an.corr_res_dif\endcsname{0.84}
\expandafter\def\csname v@an.covneg\endcsname{98.8\%}
\expandafter\def\csname v@an.covneg12\endcsname{100.0\%}
\expandafter\def\csname v@an.mean_dif\endcsname{−0.244 (t −27.67)}
\expandafter\def\csname v@an.mean_dif12\endcsname{−0.326 (t −18.54)}
\expandafter\def\csname v@an.minN\endcsname{14}
\expandafter\def\csname v@an.n.ret\endcsname{9}
\expandafter\def\csname v@an.n.ret_adj\endcsname{9}
\expandafter\def\csname v@an.p.col2\endcsname{1.219 (t 4.42)}
\expandafter\def\csname v@an.p.col3\endcsname{0.996 (t 3.75)}
\expandafter\def\csname v@an.p.col4\endcsname{0.736 (t 4.63)}
\expandafter\def\csname v@an.p.fdif\endcsname{1.261 (t 3.29)}
\expandafter\def\csname v@an.p.mean_dif\endcsname{−0.17}
\expandafter\def\csname v@an.p.sd_dif\endcsname{0.09}
\expandafter\def\csname v@an.r2_dif\endcsname{0.26}
\expandafter\def\csname v@an.r2_res\endcsname{0.62}
\expandafter\def\csname v@an.r2_res_abs\endcsname{0.61}
\expandafter\def\csname v@an.r2_res_dif\endcsname{0.87}
\expandafter\def\csname v@an.sd_dif12\endcsname{0.090}
\expandafter\def\csname v@an.sd_retadj\endcsname{29.5}
\expandafter\def\csname v@an.sens.levels\endcsname{k+10 (main) / 30 / 50}
\expandafter\def\csname v@an.sens.maxabs_t\endcsname{2.62}
\expandafter\def\csname v@an.sens.n3\endcsname{0}
\expandafter\def\csname v@an.sens.ret\endcsname{AN\_weak / LS / LS}
\expandafter\def\csname v@an.sens.ret_adj\endcsname{none / AN\_weak / none}
\expandafter\def\csname v@an.verdict.ret\endcsname{AN\_weak}
\expandafter\def\csname v@an.verdict.ret_adj\endcsname{none}
\expandafter\def\csname v@audit.TOTAL\endcsname{17}
\expandafter\def\csname v@audit.ids_in\endcsname{678}
\expandafter\def\csname v@audit.ids_out\endcsname{684}
\expandafter\def\csname v@audit.long_gap\endcsname{4}
\expandafter\def\csname v@audit.open_mismatch\endcsname{1}
\expandafter\def\csname v@audit.redenomination\endcsname{8}
\expandafter\def\csname v@audit.rows_in\endcsname{788 692}
\expandafter\def\csname v@audit.rows_out\endcsname{788 689}
\expandafter\def\csname v@audit.ticker_reuse\endcsname{2}
\expandafter\def\csname v@audit.zero_volume\endcsname{2}
\expandafter\def\csname v@base.cfg8\endcsname{3e8eb87e}
\expandafter\def\csname v@base.dir\endcsname{results/20261016T135240Z\_3e8eb87e}
\expandafter\def\csname v@base.git\endcsname{9ddffbb (dirty)}
\expandafter\def\csname v@base.rss\endcsname{865}
\expandafter\def\csname v@base.run\endcsname{20261016T135240Z\_3e8eb87e}
\expandafter\def\csname v@base.seed\endcsname{20261017}
\expandafter\def\csname v@base.t_total\endcsname{25}
\expandafter\def\csname v@d.absretMean\endcsname{4.13}
\expandafter\def\csname v@d.absretMed\endcsname{3.76}
\expandafter\def\csname v@d.auditGap\endcsname{4}
\expandafter\def\csname v@d.auditRedenom\endcsname{8}
\expandafter\def\csname v@d.auditReuse\endcsname{2}
\expandafter\def\csname v@d.avgCoins\endcsname{265}
\expandafter\def\csname v@d.cmktMax\endcsname{139.4}
\expandafter\def\csname v@d.cmktMaxMonth\endcsname{2021-03}
\expandafter\def\csname v@d.cmktMean\endcsname{4.23}
\expandafter\def\csname v@d.cmktTopA\endcsname{139.4}
\expandafter\def\csname v@d.cmktTopAC\endcsname{139.4}
\expandafter\def\csname v@d.cmktTopResA\endcsname{23.3}
\expandafter\def\csname v@d.corrLnAAbs\endcsname{0.00}
\expandafter\def\csname v@d.corrLnACAbs\endcsname{−0.15}
\expandafter\def\csname v@d.corrLnACVol\endcsname{−0.97}
\expandafter\def\csname v@d.corrResAAbs\endcsname{0.37}
\expandafter\def\csname v@d.corrVolAbs\endcsname{0.26}
\expandafter\def\csname v@d.critFmThree\endcsname{2.15}
\expandafter\def\csname v@d.critFmTwo\endcsname{2.21}
\expandafter\def\csname v@d.critHlA\endcsname{2.00}
\expandafter\def\csname v@d.critHlAC\endcsname{2.10}
\expandafter\def\csname v@d.daily_rows\endcsname{788 689}
\expandafter\def\csname v@d.daily_symbols\endcsname{684}
\expandafter\def\csname v@d.difHL\endcsname{1.21}
\expandafter\def\csname v@d.difHLT\endcsname{1.52}
\expandafter\def\csname v@d.elasMean\endcsname{1.12}
\expandafter\def\csname v@d.elasMed\endcsname{1.51}
\expandafter\def\csname v@d.exMeanA\endcsname{−0.02}
\expandafter\def\csname v@d.exMeanAC\endcsname{0.09}
\expandafter\def\csname v@d.exMeanResA\endcsname{−0.88}
\expandafter\def\csname v@d.exTA\endcsname{−0.01}
\expandafter\def\csname v@d.exTAC\endcsname{0.08}
\expandafter\def\csname v@d.exTResA\endcsname{−0.87}
\expandafter\def\csname v@d.fmAvgObs\endcsname{214}
\expandafter\def\csname v@d.fmExFourAC\endcsname{0.53}
\expandafter\def\csname v@d.fmExFourACT\endcsname{1.37}
\expandafter\def\csname v@d.fmExFourRes\endcsname{−6.65}
\expandafter\def\csname v@d.fmExFourResT\endcsname{−2.54}
\expandafter\def\csname v@d.fmExThree\endcsname{−4.55}
\expandafter\def\csname v@d.fmExThreeT\endcsname{−2.35}
\expandafter\def\csname v@d.fmExTwo\endcsname{0.33}
\expandafter\def\csname v@d.fmExTwoT\endcsname{0.97}
\expandafter\def\csname v@d.fmFourRes\endcsname{−8.10}
\expandafter\def\csname v@d.fmMonths\endcsname{87}
\expandafter\def\csname v@d.fmReproDiffFourAC\endcsname{0.000000}
\expandafter\def\csname v@d.fmReproDiffFourRes\endcsname{0.000000}
\expandafter\def\csname v@d.fmReproDiffThree\endcsname{0.000000}
\expandafter\def\csname v@d.fmReproDiffTwo\endcsname{0.000000}
\expandafter\def\csname v@d.fmTopMonthFourAC\endcsname{2021-03}
\expandafter\def\csname v@d.fmTopMonthFourRes\endcsname{2021-03}
\expandafter\def\csname v@d.fmTopMonthThree\endcsname{2021-04}
\expandafter\def\csname v@d.fmTopMonthTwo\endcsname{2021-03}
\expandafter\def\csname v@d.fmTopValFourAC\endcsname{43.20}
\expandafter\def\csname v@d.fmTopValFourRes\endcsname{−133.38}
\expandafter\def\csname v@d.fmTopValThree\endcsname{−51.45}
\expandafter\def\csname v@d.fmTopValTwo\endcsname{37.57}
\expandafter\def\csname v@d.fprFmThree\endcsname{9.5}
\expandafter\def\csname v@d.fprFmThreeHi\endcsname{14.4}
\expandafter\def\csname v@d.fprFmThreeLo\endcsname{5.8}
\expandafter\def\csname v@d.fprFmThreeNWzero\endcsname{5.0}
\expandafter\def\csname v@d.fprFmThreeNWzeroHi\endcsname{9.0}
\expandafter\def\csname v@d.fprFmThreeNWzeroLo\endcsname{2.4}
\expandafter\def\csname v@d.fprFmTwo\endcsname{8.0}
\expandafter\def\csname v@d.fprFmTwoHi\endcsname{12.7}
\expandafter\def\csname v@d.fprFmTwoLo\endcsname{4.6}
\expandafter\def\csname v@d.fprHlAC\endcsname{6.5}
\expandafter\def\csname v@d.fprHlACHi\endcsname{10.9}
\expandafter\def\csname v@d.fprHlACLo\endcsname{3.5}
\expandafter\def\csname v@d.fprThreeFmThree\endcsname{1.5}
\expandafter\def\csname v@d.fprThreeFmThreeNWzero\endcsname{1.0}
\expandafter\def\csname v@d.fprThreeFmTwo\endcsname{1.5}
\expandafter\def\csname v@d.fprThreeHlAC\endcsname{1.0}
\expandafter\def\csname v@d.hlz\endcsname{3.0}
\expandafter\def\csname v@d.janCoef\endcsname{0.70}
\expandafter\def\csname v@d.janT\endcsname{0.39}
\expandafter\def\csname v@d.leakTFmTwo\endcsname{−4.85}
\expandafter\def\csname v@d.leakTHLAC\endcsname{−6.02}
\expandafter\def\csname v@d.liqResAalpha\endcsname{−3.47}
\expandafter\def\csname v@d.liqResAalphaT\endcsname{−2.82}
\expandafter\def\csname v@d.liqResAraw\endcsname{−2.82}
\expandafter\def\csname v@d.liqResArawT\endcsname{−1.71}
\expandafter\def\csname v@d.mdePs\endcsname{2.15}
\expandafter\def\csname v@d.meanRetTopA\endcsname{140.1}
\expandafter\def\csname v@d.meanRetTopAC\endcsname{140.1}
\expandafter\def\csname v@d.meanRetTopResA\endcsname{26.4}
\expandafter\def\csname v@d.medDvolM\endcsname{3.4}
\expandafter\def\csname v@d.medianA\endcsname{−0.48}
\expandafter\def\csname v@d.medianAC\endcsname{−0.26}
\expandafter\def\csname v@d.medianResA\endcsname{−1.43}
\expandafter\def\csname v@d.nCoins\endcsname{626}
\expandafter\def\csname v@d.nMonths\endcsname{93}
\expandafter\def\csname v@d.paperAbsretMed\endcsname{1.6}
\expandafter\def\csname v@d.paperAvgObs\endcsname{1775}
\expandafter\def\csname v@d.paperMonths\endcsname{588}
\expandafter\def\csname v@d.paperSubCFour\endcsname{0.040}
\expandafter\def\csname v@d.paperSubCFourT\endcsname{0.90}
\expandafter\def\csname v@d.placeboReps\endcsname{20}
\expandafter\def\csname v@d.placeboShareMax\endcsname{10}
\expandafter\def\csname v@d.powFmTwoHalf\endcsname{95}
\expandafter\def\csname v@d.powFmTwoOne\endcsname{100}
\expandafter\def\csname v@d.powFmTwoQuarter\endcsname{35}
\expandafter\def\csname v@d.powHlACHalf\endcsname{85}
\expandafter\def\csname v@d.powHlACOne\endcsname{100}
\expandafter\def\csname v@d.powHlACQuarter\endcsname{45}
\expandafter\def\csname v@d.powReps\endcsname{20}
\expandafter\def\csname v@d.psPaper\endcsname{0.46}
\expandafter\def\csname v@d.psPaperOurSd\endcsname{0.28}
\expandafter\def\csname v@d.psUs\endcsname{1.17}
\expandafter\def\csname v@d.qResA\endcsname{5.15 / 4.49 / 4.19 / 3.72 / 4.64}
\expandafter\def\csname v@d.qratioAC\endcsname{9.4}
\expandafter\def\csname v@d.qratioAbs\endcsname{1.75}
\expandafter\def\csname v@d.rawHLA\endcsname{1.54}
\expandafter\def\csname v@d.rawHLAC\endcsname{1.70}
\expandafter\def\csname v@d.realSeHlAC\endcsname{1.52}
\expandafter\def\csname v@d.retAdjMax\endcsname{1165.7}
\expandafter\def\csname v@d.retAdjMin\endcsname{−485.9}
\expandafter\def\csname v@d.retMax\endcsname{1149.7}
\expandafter\def\csname v@d.retP1\endcsname{−54.2}
\expandafter\def\csname v@d.retP99\endcsname{158.4}
\expandafter\def\csname v@d.retSd\endcsname{43.9}
\expandafter\def\csname v@d.sdLnAC\endcsname{1.54}
\expandafter\def\csname v@d.sdPaper\endcsname{2.53}
\expandafter\def\csname v@d.seRatioHlAC\endcsname{3.4}
\expandafter\def\csname v@d.shareCovNeg\endcsname{98.8}
\expandafter\def\csname v@d.sharePosA\endcsname{47}
\expandafter\def\csname v@d.sharePosAC\endcsname{49}
\expandafter\def\csname v@d.sharePosResA\endcsname{37}
\expandafter\def\csname v@d.simAbsret\endcsname{3.49}
\expandafter\def\csname v@d.simCoins\endcsname{298}
\expandafter\def\csname v@d.simMonths\endcsname{90}
\expandafter\def\csname v@d.simReps\endcsname{200}
\expandafter\def\csname v@d.simSeHlAC\endcsname{0.45}
\expandafter\def\csname v@d.subAFm\endcsname{1.98}
\expandafter\def\csname v@d.subAFmT\endcsname{2.59}
\expandafter\def\csname v@d.subAHLAC\endcsname{3.39}
\expandafter\def\csname v@d.subAHLACt\endcsname{1.21}
\expandafter\def\csname v@d.subAObs\endcsname{102}
\expandafter\def\csname v@d.subARange\endcsname{2019-01..2022-10}
\expandafter\def\csname v@d.subBFm\endcsname{−0.28}
\expandafter\def\csname v@d.subBFmT\endcsname{−1.09}
\expandafter\def\csname v@d.subBHLAC\endcsname{0.05}
\expandafter\def\csname v@d.subBHLACt\endcsname{0.05}
\expandafter\def\csname v@d.subBObs\endcsname{309}
\expandafter\def\csname v@d.subBRange\endcsname{2022-11..2026-09}
\expandafter\def\csname v@d.topMonthA\endcsname{2021-03}
\expandafter\def\csname v@d.topMonthAC\endcsname{2021-03}
\expandafter\def\csname v@d.topMonthResA\endcsname{2019-03}
\expandafter\def\csname v@d.topValA\endcsname{145.0}
\expandafter\def\csname v@d.topValAC\endcsname{150.3}
\expandafter\def\csname v@d.topValResA\endcsname{33.4}
\expandafter\def\csname v@d.vwResAalpha\endcsname{−3.29}
\expandafter\def\csname v@d.vwResAalphaT\endcsname{−2.38}
\expandafter\def\csname v@d.vwResAraw\endcsname{−2.80}
\expandafter\def\csname v@d.vwResArawT\endcsname{−1.82}
\expandafter\def\csname v@d.winIvolT\endcsname{−3.15}
\expandafter\def\csname v@ed.alone_pos\endcsname{5 of 6}
\expandafter\def\csname v@ed.bidask_diff\endcsname{0.0}
\expandafter\def\csname v@ed.bidask_n\endcsname{300}
\expandafter\def\csname v@ed.breakeven\endcsname{1.52\%}
\expandafter\def\csname v@ed.corr.A\endcsname{0.181}
\expandafter\def\csname v@ed.corr.A.paper\endcsname{0.731}
\expandafter\def\csname v@ed.corr.A_C\endcsname{0.158}
\expandafter\def\csname v@ed.corr.A_C.paper\endcsname{0.705}
\expandafter\def\csname v@ed.corr.ResA\endcsname{0.109}
\expandafter\def\csname v@ed.corr.ResA.paper\endcsname{0.215}
\expandafter\def\csname v@ed.cost_pos\endcsname{5 of 6}
\expandafter\def\csname v@ed.fee\endcsname{0.1\%}
\expandafter\def\csname v@ed.fm1\endcsname{−0.017 (t −0.04)}
\expandafter\def\csname v@ed.fm1.paper\endcsname{−0.248 (t −1.94)}
\expandafter\def\csname v@ed.fm2ac\endcsname{0.762 (t 1.51)}
\expandafter\def\csname v@ed.gap_bp\endcsname{3.3}
\expandafter\def\csname v@ed.gap_zero\endcsname{35\%}
\expandafter\def\csname v@ed.gross\endcsname{1.65 (t 1.07)}
\expandafter\def\csname v@ed.hs0_long\endcsname{42\%}
\expandafter\def\csname v@ed.hs0_short\endcsname{54\%}
\expandafter\def\csname v@ed.maxabs_t\endcsname{1.94}
\expandafter\def\csname v@ed.months\endcsname{92}
\expandafter\def\csname v@ed.n\endcsname{26 046}
\expandafter\def\csname v@ed.net\endcsname{1.57 (t 1.02)}
\expandafter\def\csname v@ed.netfee\endcsname{1.46 (t 0.95)}
\expandafter\def\csname v@ed.nonpos\endcsname{50.7\%}
\expandafter\def\csname v@ed.t10cost\endcsname{−10.11 (t −1.07)}
\expandafter\def\csname v@ed.t10non\endcsname{0.547 (t 1.17)}
\expandafter\def\csname v@ed.turn_long\endcsname{0.62}
\expandafter\def\csname v@ed.turn_short\endcsname{0.46}
\expandafter\def\csname v@ed.w2\endcsname{0.9996}
\expandafter\def\csname v@f.cmkt_mean\endcsname{4.23}
\expandafter\def\csname v@git.firstrun\endcsname{20261016T134056Z\_18099336}
\expandafter\def\csname v@git.firstrun.cfg8\endcsname{18099336}
\expandafter\def\csname v@git.firstrun.sha\endcsname{8b60380 (dirty)}
\expandafter\def\csname v@git.prereg\endcsname{a90a72b}
\expandafter\def\csname v@git.prereg.cfgmatch\endcsname{совпадает}
\expandafter\def\csname v@git.refine\endcsname{8b60380}
\expandafter\def\csname v@hf.b0\endcsname{−0.435 (t −1.56)}
\expandafter\def\csname v@hf.b1\endcsname{−2.59 (t −4.03)}
\expandafter\def\csname v@hf.b1.paper\endcsname{−0.034 (t −0.67)}
\expandafter\def\csname v@hf.b2ac\endcsname{1.06 (t 1.97)}
\expandafter\def\csname v@hf.b2ac.paper\endcsname{0.18 (t 3.80)}
\expandafter\def\csname v@hf.b2lam\endcsname{−4.30 (t −3.57)}
\expandafter\def\csname v@hf.b2lam.paper\endcsname{−0.155 (t −2.51)}
\expandafter\def\csname v@hf.coinmonths\endcsname{13 708}
\expandafter\def\csname v@hf.corr.A\endcsname{0.615}
\expandafter\def\csname v@hf.corr.A.paper\endcsname{0.737}
\expandafter\def\csname v@hf.corr.A_C\endcsname{0.547}
\expandafter\def\csname v@hf.corr.A_C.paper\endcsname{0.747}
\expandafter\def\csname v@hf.corr.ResA\endcsname{0.294}
\expandafter\def\csname v@hf.corr.ResA.paper\endcsname{0.149}
\expandafter\def\csname v@hf.first\endcsname{2023-10}
\expandafter\def\csname v@hf.gb5m\endcsname{4.3}
\expandafter\def\csname v@hf.gbagg\endcsname{1.0}
\expandafter\def\csname v@hf.h1.ac_t\endcsname{4.07}
\expandafter\def\csname v@hf.h2.ac_t\endcsname{0.13}
\expandafter\def\csname v@hf.lam_tmax\endcsname{5.15}
\expandafter\def\csname v@hf.lam_tmin\endcsname{2.34}
\expandafter\def\csname v@hf.last\endcsname{2026-09}
\expandafter\def\csname v@hf.months\endcsname{36}
\expandafter\def\csname v@hf.obs\endcsname{318}
\expandafter\def\csname v@hf.t10.c1\endcsname{−1.43 (t −4.03)}
\expandafter\def\csname v@hf.t10.c1.paper\endcsname{−0.047 (t −1.11)}
\expandafter\def\csname v@hf.t10.c3\endcsname{−1.42 (t −3.92)}
\expandafter\def\csname v@hf.t10.c3.paper\endcsname{0.015 (t 0.35)}
\expandafter\def\csname v@hf.t10.n2\endcsname{1.04 (t 2.11)}
\expandafter\def\csname v@hf.t10.n2.paper\endcsname{0.186 (t 3.04)}
\expandafter\def\csname v@hf.t10.n3\endcsname{0.957 (t 1.88)}
\expandafter\def\csname v@hf.t10.n3.paper\endcsname{0.191 (t 3.10)}
\expandafter\def\csname v@hf.t10r2\endcsname{0.68}
\expandafter\def\csname v@hf.val.corr\endcsname{0.86}
\expandafter\def\csname v@hf.val.n\endcsname{69}
\expandafter\def\csname v@hf.vol.b1\endcsname{−3.92 (t −3.78)}
\expandafter\def\csname v@hf.vol.b2ac\endcsname{0.074 (t 0.08)}
\expandafter\def\csname v@il.fm\endcsname{−0.441 (t −2.19)}
\expandafter\def\csname v@il.long\endcsname{−1.66 (t −2.59)}
\expandafter\def\csname v@il.long.paper\endcsname{−0.016 (t −1.11)}
\expandafter\def\csname v@il.ls\endcsname{−4.30 (t −3.58)}
\expandafter\def\csname v@il.ls.paper\endcsname{−0.048 (t −2.38)}
\expandafter\def\csname v@il.n\endcsname{88}
\expandafter\def\csname v@il.post.long\endcsname{−0.858 (t −1.12)}
\expandafter\def\csname v@il.post.ls\endcsname{−3.16 (t −2.60)}
\expandafter\def\csname v@il.post.short\endcsname{2.31 (t 3.13)}
\expandafter\def\csname v@il.short\endcsname{2.64 (t 3.78)}
\expandafter\def\csname v@il.short.paper\endcsname{0.031 (t 2.18)}
\expandafter\def\csname v@il.split\endcsname{−5.60 (t −1.45)}
\expandafter\def\csname v@il.split.paper\endcsname{−0.02 (t −0.08)}
\expandafter\def\csname v@il.var_196\endcsname{5}
\expandafter\def\csname v@il.var_3\endcsname{3}
\expandafter\def\csname v@il.var_n\endcsname{6}
\expandafter\def\csname v@il.var_tmax\endcsname{−1.46}
\expandafter\def\csname v@il.var_tmin\endcsname{−5.51}
\expandafter\def\csname v@lg.all.196\endcsname{191}
\expandafter\def\csname v@lg.all.3\endcsname{66}
\expandafter\def\csname v@lg.all.3_by_source\endcsname{amihud\_noh 1, hf\_lambda 25, moex 29, openap\_us 8, sentiment 3}
\expandafter\def\csname v@lg.all.bh\endcsname{90}
\expandafter\def\csname v@lg.all.bonf\endcsname{14}
\expandafter\def\csname v@lg.all.crit\endcsname{3.99}
\expandafter\def\csname v@lg.all.exp\endcsname{37.5}
\expandafter\def\csname v@lg.all.n\endcsname{750}
\expandafter\def\csname v@lg.base.196\endcsname{16}
\expandafter\def\csname v@lg.base.3\endcsname{0}
\expandafter\def\csname v@lg.base.bh\endcsname{0}
\expandafter\def\csname v@lg.base.bonf\endcsname{0}
\expandafter\def\csname v@lg.base.crit\endcsname{3.40}
\expandafter\def\csname v@lg.base.exp\endcsname{3.6}
\expandafter\def\csname v@lg.base.n\endcsname{73}
\expandafter\def\csname v@loyo.fm2\endcsname{0.96..1.94}
\expandafter\def\csname v@loyo.fm2.n196\endcsname{0}
\expandafter\def\csname v@loyo.fm2.years\endcsname{8}
\expandafter\def\csname v@loyo.hlAC\endcsname{0.42..1.32}
\expandafter\def\csname v@loyo.hlAC.n196\endcsname{0}
\expandafter\def\csname v@loyo.hlAC.years\endcsname{8}
\expandafter\def\csname v@mde.fm2\endcsname{2.15}
\expandafter\def\csname v@mde.fm2.paper\endcsname{0.46}
\expandafter\def\csname v@mde.fm2.paper_oursd\endcsname{0.28}
\expandafter\def\csname v@mde.fm2.paper_sd\endcsname{2.53}
\expandafter\def\csname v@mde.fm2.paper_slope\endcsname{0.183}
\expandafter\def\csname v@mde.fm2.ratio\endcsname{4.6}
\expandafter\def\csname v@mde.fm2.sd\endcsname{1.54}
\expandafter\def\csname v@mde.fm2.sim\endcsname{0.44}
\expandafter\def\csname v@mde.fm2.slope\endcsname{1.39}
\expandafter\def\csname v@mde.fm2.slope_ratio\endcsname{7.6}
\expandafter\def\csname v@mde.hlA\endcsname{4.42}
\expandafter\def\csname v@mde.hlA.paper\endcsname{0.56}
\expandafter\def\csname v@mde.hlA.ratio\endcsname{7.9}
\expandafter\def\csname v@mde.hlAC\endcsname{4.27}
\expandafter\def\csname v@mde.hlAC.paper\endcsname{0.61}
\expandafter\def\csname v@mde.hlAC.ratio\endcsname{7.0}
\expandafter\def\csname v@mx.alpha_months\endcsname{138}
\expandafter\def\csname v@mx.avg\endcsname{164}
\expandafter\def\csname v@mx.corrAAC\endcsname{0.891}
\expandafter\def\csname v@mx.corrAAC.paper\endcsname{0.899}
\expandafter\def\csname v@mx.corrAT\endcsname{0.910}
\expandafter\def\csname v@mx.corrAT.paper\endcsname{0.746}
\expandafter\def\csname v@mx.corrAabs\endcsname{0.187}
\expandafter\def\csname v@mx.corrAabs.paper\endcsname{0.489}
\expandafter\def\csname v@mx.ex_topyear\endcsname{0.424 (t 0.97)}
\expandafter\def\csname v@mx.first\endcsname{2014-08}
\expandafter\def\csname v@mx.fm1\endcsname{0.313 (t 3.01)}
\expandafter\def\csname v@mx.fm1.paper\endcsname{0.119 (t 3.01)}
\expandafter\def\csname v@mx.fm2\endcsname{0.308 (t 2.87)}
\expandafter\def\csname v@mx.fm2.paper\endcsname{0.183 (t 4.79)}
\expandafter\def\csname v@mx.fm4_ex_topyear\endcsname{0.226 (t 2.79)}
\expandafter\def\csname v@mx.fm4_loyo_tmin\endcsname{2.34}
\expandafter\def\csname v@mx.fm4ac\endcsname{0.298 (t 2.85)}
\expandafter\def\csname v@mx.fm4ac.paper\endcsname{0.13 (t 3.31)}
\expandafter\def\csname v@mx.fm4res\endcsname{−0.038 (t −0.13)}
\expandafter\def\csname v@mx.fm4res.paper\endcsname{−0.248 (t −3.14)}
\expandafter\def\csname v@mx.fm5ac\endcsname{0.326 (t 3.19)}
\expandafter\def\csname v@mx.fm5ac.paper\endcsname{0.12 (t 3.11)}
\expandafter\def\csname v@mx.fmmonths\endcsname{121}
\expandafter\def\csname v@mx.hlA\endcsname{1.31 (t 1.81)}
\expandafter\def\csname v@mx.hlA.paper\endcsname{0.56 (t 2.36)}
\expandafter\def\csname v@mx.hlAC\endcsname{1.11 (t 1.53)}
\expandafter\def\csname v@mx.hlAC.paper\endcsname{0.61 (t 2.95)}
\expandafter\def\csname v@mx.hlACalpha\endcsname{0.034 (t 0.13)}
\expandafter\def\csname v@mx.hlA_same\endcsname{1.50 (t 2.09)}
\expandafter\def\csname v@mx.hlz_list\endcsname{FM (1) lnA t 3.01; FM (5) lnA\_C t 3.19; FM (5) lnAT\_C t 3.04}
\expandafter\def\csname v@mx.hlz_n\endcsname{3}
\expandafter\def\csname v@mx.last\endcsname{2026-09}
\expandafter\def\csname v@mx.lnme_q1\endcsname{13.22}
\expandafter\def\csname v@mx.lnme_q5\endcsname{8.49}
\expandafter\def\csname v@mx.months\endcsname{145}
\expandafter\def\csname v@mx.smb\endcsname{0.906 (t 8.54)}
\expandafter\def\csname v@mx.stocks\endcsname{306}
\expandafter\def\csname v@mx.topyear\endcsname{2023}
\expandafter\def\csname v@mx.topyear_share\endcsname{65\%}
\expandafter\def\csname v@mx.vol_196\endcsname{8}
\expandafter\def\csname v@mx.vol_n\endcsname{8}
\expandafter\def\csname v@n.avgcoins\endcsname{265}
\expandafter\def\csname v@n.avgcoins_fm_all\endcsname{201}
\expandafter\def\csname v@n.coinmonths\endcsname{24 628}
\expandafter\def\csname v@n.coins\endcsname{626}
\expandafter\def\csname v@n.first\endcsname{2019-01}
\expandafter\def\csname v@n.fm_dropped\endcsname{6}
\expandafter\def\csname v@n.fmmonths\endcsname{87}
\expandafter\def\csname v@n.fmobs\endcsname{214}
\expandafter\def\csname v@n.last\endcsname{2026-09}
\expandafter\def\csname v@n.months\endcsname{93}
\expandafter\def\csname v@op.corr.full\endcsname{0.67}
\expandafter\def\csname v@op.corr.post\endcsname{0.49}
\expandafter\def\csname v@op.corr.pre\endcsname{0.65}
\expandafter\def\csname v@op.dv.decay\endcsname{−0.732 (t −3.12)}
\expandafter\def\csname v@op.dv.decay_lags\endcsname{0 / 3 / 6 / 12}
\expandafter\def\csname v@op.dv.decay_ts\endcsname{−2.64 / −3.00 / −3.12 / −3.11}
\expandafter\def\csname v@op.dv.post\endcsname{−0.013 (t −0.09)}
\expandafter\def\csname v@op.dv.pre\endcsname{0.719 (t 3.98)}
\expandafter\def\csname v@op.ill.decay\endcsname{−0.466 (t −3.08)}
\expandafter\def\csname v@op.ill.decay_lags\endcsname{0 / 3 / 6 / 12}
\expandafter\def\csname v@op.ill.decay_ts\endcsname{−2.58 / −3.00 / −3.08 / −2.97}
\expandafter\def\csname v@op.ill.post\endcsname{−0.142 (t −1.47)}
\expandafter\def\csname v@op.ill.pre\endcsname{0.324 (t 2.78)}
\expandafter\def\csname v@op.paper.ff4\endcsname{0.35 (t 2.31)}
\expandafter\def\csname v@op.paper.five\endcsname{−0.08 (t −1.66)}
\expandafter\def\csname v@op.paper.one\endcsname{−0.04 (t −0.77)}
\expandafter\def\csname v@op.paper.raw\endcsname{0.56 (t 2.36)}
\expandafter\def\csname v@op.post.ff4\endcsname{0.165 (t 1.66)}
\expandafter\def\csname v@op.post.five\endcsname{0.077 (t 0.80)}
\expandafter\def\csname v@op.post.one\endcsname{−0.137 (t −1.40)}
\expandafter\def\csname v@op.post_range\endcsname{2013-01..2024-12}
\expandafter\def\csname v@op.pre.ff4\endcsname{0.136 (t 1.60)}
\expandafter\def\csname v@op.pre.five\endcsname{−0.018 (t −0.25)}
\expandafter\def\csname v@op.pre.one\endcsname{0.054 (t 0.63)}
\expandafter\def\csname v@plc.leak_maxabs_t\endcsname{6.02}
\expandafter\def\csname v@plc.leak_minabs_t\endcsname{2.58}
\expandafter\def\csname v@plc.shuffle_maxabs_meant\endcsname{0.32}
\expandafter\def\csname v@plc.shuffle_reps\endcsname{20}
\expandafter\def\csname v@plc.shuffle_sd_t_max\endcsname{1.12}
\expandafter\def\csname v@plc.shuffle_sd_t_min\endcsname{0.77}
\expandafter\def\csname v@plc.shuffle_share196_max\endcsname{10\%}
\expandafter\def\csname v@pow.0.00.fm2\endcsname{10\%}
\expandafter\def\csname v@pow.0.00.hlAC\endcsname{10\%}
\expandafter\def\csname v@pow.0.25.fm2\endcsname{35\%}
\expandafter\def\csname v@pow.0.25.hlAC\endcsname{45\%}
\expandafter\def\csname v@pow.0.50.fm2\endcsname{95\%}
\expandafter\def\csname v@pow.0.50.hlAC\endcsname{85\%}
\expandafter\def\csname v@pow.1.00.fm2\endcsname{100\%}
\expandafter\def\csname v@pow.1.00.hlAC\endcsname{100\%}
\expandafter\def\csname v@pow.reps\endcsname{20}
\expandafter\def\csname v@pvu.ANDIF\endcsname{−1.64 (t −0.44)}
\expandafter\def\csname v@pvu.ANDIF.match\endcsname{no}
\expandafter\def\csname v@pvu.ANDIF.paper\endcsname{0.996 (t 3.75)}
\expandafter\def\csname v@pvu.ANResLnA\endcsname{−6.33 (t −1.98)}
\expandafter\def\csname v@pvu.ANResLnA.match\endcsname{no}
\expandafter\def\csname v@pvu.ANResLnA.paper\endcsname{0.736 (t 4.63)}
\expandafter\def\csname v@pvu.FMResLnA\endcsname{−5.09 (t −2.46)}
\expandafter\def\csname v@pvu.FMResLnA.match\endcsname{yes}
\expandafter\def\csname v@pvu.FMResLnA.paper\endcsname{−0.383 (t −5.07)}
\expandafter\def\csname v@pvu.FMResLnAfour\endcsname{−8.10 (t −2.37)}
\expandafter\def\csname v@pvu.FMResLnAfour.match\endcsname{yes}
\expandafter\def\csname v@pvu.FMResLnAfour.paper\endcsname{−0.248 (t −3.14)}
\expandafter\def\csname v@pvu.FMResLnAivol\endcsname{−6.33 (t −1.98)}
\expandafter\def\csname v@pvu.FMResLnAivol.match\endcsname{yes}
\expandafter\def\csname v@pvu.FMResLnAivol.paper\endcsname{−0.303 (t −4.65)}
\expandafter\def\csname v@pvu.FMivol\endcsname{−101.34 (t −2.64)}
\expandafter\def\csname v@pvu.FMivol.match\endcsname{sign only}
\expandafter\def\csname v@pvu.FMivol.paper\endcsname{−2.903 (t −0.57)}
\expandafter\def\csname v@pvu.FMlnA\endcsname{0.580 (t 1.24)}
\expandafter\def\csname v@pvu.FMlnA.match\endcsname{sign only}
\expandafter\def\csname v@pvu.FMlnA.paper\endcsname{0.119 (t 3.01)}
\expandafter\def\csname v@pvu.FMlnAC\endcsname{0.760 (t 1.53)}
\expandafter\def\csname v@pvu.FMlnAC.match\endcsname{sign only}
\expandafter\def\csname v@pvu.FMlnAC.paper\endcsname{0.183 (t 4.79)}
\expandafter\def\csname v@pvu.FMlnACfive\endcsname{1.19 (t 1.70)}
\expandafter\def\csname v@pvu.FMlnACfive.match\endcsname{sign only}
\expandafter\def\csname v@pvu.FMlnACfive.paper\endcsname{0.12 (t 3.11)}
\expandafter\def\csname v@pvu.FMlnACfour\endcsname{1.02 (t 1.67)}
\expandafter\def\csname v@pvu.FMlnACfour.match\endcsname{sign only}
\expandafter\def\csname v@pvu.FMlnACfour.paper\endcsname{0.13 (t 3.31)}
\expandafter\def\csname v@pvu.HLACalphaCoef\endcsname{0.781 (t 0.62)}
\expandafter\def\csname v@pvu.HLACalphaCoef.match\endcsname{sign only}
\expandafter\def\csname v@pvu.HLACalphaCoef.paper\endcsname{0.44 (t 3.20)}
\expandafter\def\csname v@pvu.HLACrawCoef\endcsname{1.70 (t 1.12)}
\expandafter\def\csname v@pvu.HLACrawCoef.match\endcsname{sign only}
\expandafter\def\csname v@pvu.HLACrawCoef.paper\endcsname{0.61 (t 2.95)}
\expandafter\def\csname v@pvu.HLAalphaCoef\endcsname{0.556 (t 0.48)}
\expandafter\def\csname v@pvu.HLAalphaCoef.match\endcsname{sign only}
\expandafter\def\csname v@pvu.HLAalphaCoef.paper\endcsname{0.35 (t 2.31)}
\expandafter\def\csname v@pvu.HLArawCoef\endcsname{1.54 (t 0.98)}
\expandafter\def\csname v@pvu.HLArawCoef.match\endcsname{sign only}
\expandafter\def\csname v@pvu.HLArawCoef.paper\endcsname{0.56 (t 2.36)}
\expandafter\def\csname v@pvu.HLResAalphaCoef\endcsname{−1.00 (t −1.01)}
\expandafter\def\csname v@pvu.HLResAalphaCoef.match\endcsname{yes}
\expandafter\def\csname v@pvu.HLResAalphaCoef.paper\endcsname{−0.16 (t −0.96)}
\expandafter\def\csname v@pvu.HLResArawCoef\endcsname{−0.510 (t −0.49)}
\expandafter\def\csname v@pvu.HLResArawCoef.match\endcsname{yes}
\expandafter\def\csname v@pvu.HLResArawCoef.paper\endcsname{−0.17 (t −1.05)}
\expandafter\def\csname v@pvu.IMLAalpha\endcsname{0.234 (t 0.50)}
\expandafter\def\csname v@pvu.IMLAalpha.match\endcsname{yes}
\expandafter\def\csname v@pvu.IMLAalpha.paper\endcsname{−0.04 (t −0.77)}
\expandafter\def\csname v@pvu.IMLAalphaFive\endcsname{0.155 (t 0.35)}
\expandafter\def\csname v@pvu.IMLAalphaFive.match\endcsname{yes}
\expandafter\def\csname v@pvu.IMLAalphaFive.paper\endcsname{−0.08 (t −1.66)}
\expandafter\def\csname v@pvu.QratioAC\endcsname{9.40}
\expandafter\def\csname v@pvu.QratioAC.match\endcsname{no}
\expandafter\def\csname v@pvu.QratioAC.paper\endcsname{116.60}
\expandafter\def\csname v@pvu.QratioAbs\endcsname{1.75}
\expandafter\def\csname v@pvu.QratioAbs.match\endcsname{yes}
\expandafter\def\csname v@pvu.QratioAbs.paper\endcsname{2.18}
\expandafter\def\csname v@pvu.RawResLnA\endcsname{−3.70 (t −2.11)}
\expandafter\def\csname v@pvu.RawResLnA.match\endcsname{sign only}
\expandafter\def\csname v@pvu.RawResLnA.paper\endcsname{−0.143 (t −1.13)}
\expandafter\def\csname v@pvu.RawResLnAfive\endcsname{−6.04 (t −2.62)}
\expandafter\def\csname v@pvu.RawResLnAfive.match\endcsname{yes}
\expandafter\def\csname v@pvu.RawResLnAfive.paper\endcsname{−0.185 (t −2.23)}
\expandafter\def\csname v@pvu.RawlnA\endcsname{0.688 (t 1.39)}
\expandafter\def\csname v@pvu.RawlnA.match\endcsname{sign only}
\expandafter\def\csname v@pvu.RawlnA.paper\endcsname{0.081 (t 2.08)}
\expandafter\def\csname v@pvu.RawlnAC\endcsname{0.860 (t 1.70)}
\expandafter\def\csname v@pvu.RawlnAC.match\endcsname{yes}
\expandafter\def\csname v@pvu.RawlnAC.paper\endcsname{0.095 (t 1.93)}
\expandafter\def\csname v@pvu.SubAResLnA\endcsname{−13.90 (t −2.29)}
\expandafter\def\csname v@pvu.SubAResLnA.match\endcsname{yes}
\expandafter\def\csname v@pvu.SubAResLnA.paper\endcsname{−0.222 (t −2.24)}
\expandafter\def\csname v@pvu.SubAlnAC\endcsname{1.98 (t 2.59)}
\expandafter\def\csname v@pvu.SubAlnAC.match\endcsname{yes}
\expandafter\def\csname v@pvu.SubAlnAC.paper\endcsname{0.272 (t 4.71)}
\expandafter\def\csname v@pvu.SubBResLnA\endcsname{−3.17 (t −1.40)}
\expandafter\def\csname v@pvu.SubBResLnA.match\endcsname{sign only}
\expandafter\def\csname v@pvu.SubBResLnA.paper\endcsname{−0.275 (t −2.21)}
\expandafter\def\csname v@pvu.SubBlnAC\endcsname{−0.280 (t −1.09)}
\expandafter\def\csname v@pvu.SubBlnAC.match\endcsname{no}
\expandafter\def\csname v@pvu.SubBlnAC.paper\endcsname{0.902 (t 1.98)}
\expandafter\def\csname v@pvu.TsixLnVOL\endcsname{−1.31 (t −1.81)}
\expandafter\def\csname v@pvu.TsixLnVOL.match\endcsname{sign only}
\expandafter\def\csname v@pvu.TsixLnVOL.paper\endcsname{−0.158 (t −3.65)}
\expandafter\def\csname v@pvu.TsixResLnA\endcsname{−8.76 (t −2.29)}
\expandafter\def\csname v@pvu.TsixResLnA.match\endcsname{yes}
\expandafter\def\csname v@pvu.TsixResLnA.paper\endcsname{−0.167 (t −2.74)}
\expandafter\def\csname v@pvu.TtwolnAC\endcsname{0.814 (t 1.46)}
\expandafter\def\csname v@pvu.TtwolnAC.match\endcsname{sign only}
\expandafter\def\csname v@pvu.TtwolnAC.paper\endcsname{0.165 (t 3.92)}
\expandafter\def\csname v@pvu.TtwolnAbs\endcsname{−9.05 (t −2.50)}
\expandafter\def\csname v@pvu.TtwolnAbs.match\endcsname{yes}
\expandafter\def\csname v@pvu.TtwolnAbs.paper\endcsname{−0.319 (t −2.13)}
\expandafter\def\csname v@pvu.corrAAC\endcsname{0.956}
\expandafter\def\csname v@pvu.corrAAC.match\endcsname{yes}
\expandafter\def\csname v@pvu.corrAAC.paper\endcsname{0.899}
\expandafter\def\csname v@pvu.corrACabs\endcsname{−0.247}
\expandafter\def\csname v@pvu.corrACabs.match\endcsname{no}
\expandafter\def\csname v@pvu.corrACabs.paper\endcsname{0.394}
\expandafter\def\csname v@pvu.corrAabs\endcsname{−0.133}
\expandafter\def\csname v@pvu.corrAabs.match\endcsname{no}
\expandafter\def\csname v@pvu.corrAabs.paper\endcsname{0.489}
\expandafter\def\csname v@q.A\endcsname{3.40 / 4.55 / 4.93 / 4.24 / 4.94}
\expandafter\def\csname v@q.A.max\endcsname{4.9}
\expandafter\def\csname v@q.A.min\endcsname{3.4}
\expandafter\def\csname v@q.A.paper\endcsname{0.96 / 1.16 / 1.23 / 1.29 / 1.53}
\expandafter\def\csname v@q.A_C\endcsname{3.41 / 3.89 / 5.38 / 4.22 / 5.11}
\expandafter\def\csname v@q.A_C.max\endcsname{5.4}
\expandafter\def\csname v@q.A_C.min\endcsname{3.4}
\expandafter\def\csname v@q.A_C.paper\endcsname{0.96 / 1.13 / 1.23 / 1.28 / 1.57}
\expandafter\def\csname v@q.ResA\endcsname{5.15 / 4.49 / 4.19 / 3.72 / 4.64}
\expandafter\def\csname v@q.ResA.max\endcsname{5.2}
\expandafter\def\csname v@q.ResA.min\endcsname{3.7}
\expandafter\def\csname v@q.ResA.paper\endcsname{1.39 / 1.29 / 1.19 / 1.10 / 1.21}
\expandafter\def\csname v@rb.block\endcsname{6}
\expandafter\def\csname v@rb.ci.fm1_lnA\endcsname{[−0.22; 1.62]}
\expandafter\def\csname v@rb.ci.fm2_lnA_C\endcsname{[−0.10; 1.87]}
\expandafter\def\csname v@rb.ci.fm3_ResLnA\endcsname{[−8.98; −1.52]}
\expandafter\def\csname v@rb.ci.fm4_ResLnA\endcsname{[−15.29; −2.45]}
\expandafter\def\csname v@rb.ci.fm4_lnA_C\endcsname{[0.00; 2.39]}
\expandafter\def\csname v@rb.ci.hl_A\endcsname{[−1.41; 5.33]}
\expandafter\def\csname v@rb.ci.hl_A_C\endcsname{[−1.24; 5.33]}
\expandafter\def\csname v@rb.ci.hl_A_C_alpha\endcsname{[−2.12; 2.67]}
\expandafter\def\csname v@rb.ci.hl_A_alpha\endcsname{[−2.19; 2.35]}
\expandafter\def\csname v@rb.ci.hl_ResA\endcsname{[−2.76; 1.20]}
\expandafter\def\csname v@rb.ci.hl_ResA_alpha\endcsname{[−3.13; 0.62]}
\expandafter\def\csname v@rb.draws\endcsname{2000}
\expandafter\def\csname v@rb.excl\endcsname{3}
\expandafter\def\csname v@rb.n\endcsname{11}
\expandafter\def\csname v@rb.robust\endcsname{2}
\expandafter\def\csname v@rb.seeds\endcsname{20}
\expandafter\def\csname v@rb.seedshare.fm1_lnA\endcsname{0/20}
\expandafter\def\csname v@rb.seedshare.fm2_lnA_C\endcsname{0/20}
\expandafter\def\csname v@rb.seedshare.fm3_ResLnA\endcsname{20/20}
\expandafter\def\csname v@rb.seedshare.fm4_ResLnA\endcsname{20/20}
\expandafter\def\csname v@rb.seedshare.fm4_lnA_C\endcsname{11/20}
\expandafter\def\csname v@rb.seedshare.hl_A\endcsname{0/20}
\expandafter\def\csname v@rb.seedshare.hl_A_C\endcsname{0/20}
\expandafter\def\csname v@rb.seedshare.hl_A_C_alpha\endcsname{0/20}
\expandafter\def\csname v@rb.seedshare.hl_A_alpha\endcsname{0/20}
\expandafter\def\csname v@rb.seedshare.hl_ResA\endcsname{0/20}
\expandafter\def\csname v@rb.seedshare.hl_ResA_alpha\endcsname{0/20}
\expandafter\def\csname v@rob.n\endcsname{45}
\expandafter\def\csname v@rob.same\endcsname{31}
\expandafter\def\csname v@run.amihud_noh\endcsname{20261016T141015Z\_amihud\_noh}
\expandafter\def\csname v@run.edge_spread\endcsname{20261016T141008Z\_3e8eb87e\_edge\_spread}
\expandafter\def\csname v@run.hf_lambda\endcsname{hf\_lambda\_20261016T154107Z\_3e8eb87e}
\expandafter\def\csname v@run.moex\endcsname{moex\_20261016T144811Z\_b4264d55}
\expandafter\def\csname v@run.openap_us\endcsname{20261016T140227Z\_openap\_us}
\expandafter\def\csname v@run.robustness\endcsname{robustness\_20261016T162250Z}
\expandafter\def\csname v@run.sentiment\endcsname{sentiment\_20261016T140147Z}
\expandafter\def\csname v@rv.liq.hlAC\endcsname{0.157 (t 0.07)}
\expandafter\def\csname v@sc.an_n\endcsname{2}
\expandafter\def\csname v@sc.an_opp\endcsname{2}
\expandafter\def\csname v@sc.hlz\endcsname{0}
\expandafter\def\csname v@sc.ls_n\endcsname{28}
\expandafter\def\csname v@sc.lvl_n\endcsname{5}
\expandafter\def\csname v@sc.lvl_yes\endcsname{2}
\expandafter\def\csname v@sc.no\endcsname{1}
\expandafter\def\csname v@sc.res_n\endcsname{8}
\expandafter\def\csname v@sc.res_sig\endcsname{7}
\expandafter\def\csname v@sc.samesign\endcsname{25}
\expandafter\def\csname v@sc.sig\endcsname{10}
\expandafter\def\csname v@sc.sign\endcsname{14}
\expandafter\def\csname v@sc.vol_n\endcsname{11}
\expandafter\def\csname v@sc.vol_sig\endcsname{0}
\expandafter\def\csname v@sc.yes\endcsname{13}
\expandafter\def\csname v@se.fm\endcsname{0.586 (t 1.99)}
\expandafter\def\csname v@se.long\endcsname{1.20 (t 1.41)}
\expandafter\def\csname v@se.long.paper\endcsname{−0.01 (t −0.10)}
\expandafter\def\csname v@se.ls\endcsname{2.57 (t 1.84)}
\expandafter\def\csname v@se.ls.paper\endcsname{0.24 (t 1.81)}
\expandafter\def\csname v@se.n\endcsname{93}
\expandafter\def\csname v@se.post.long\endcsname{−0.120 (t −0.11)}
\expandafter\def\csname v@se.post.ls\endcsname{0.604 (t 0.32)}
\expandafter\def\csname v@se.post.short\endcsname{−0.724 (t −0.63)}
\expandafter\def\csname v@se.short\endcsname{−1.37 (t −1.70)}
\expandafter\def\csname v@se.short.paper\endcsname{−0.25 (t −2.67)}
\expandafter\def\csname v@se.split\endcsname{3.64 (t 1.21)}
\expandafter\def\csname v@se.split.paper\endcsname{0.36 (t 1.08)}
\expandafter\def\csname v@se.var_196\endcsname{0}
\expandafter\def\csname v@se.var_3\endcsname{0}
\expandafter\def\csname v@se.var_n\endcsname{4}
\expandafter\def\csname v@se.var_tmax\endcsname{1.84}
\expandafter\def\csname v@se.var_tmin\endcsname{0.12}
\expandafter\def\csname v@sim.0.fm2\endcsname{8.0\%}
\expandafter\def\csname v@sim.0.fm2.3\endcsname{1.5\%}
\expandafter\def\csname v@sim.0.fm2nw0\endcsname{6.0\%}
\expandafter\def\csname v@sim.0.fm2nw0.3\endcsname{0.0\%}
\expandafter\def\csname v@sim.0.fm3\endcsname{9.5\%}
\expandafter\def\csname v@sim.0.fm3.3\endcsname{1.5\%}
\expandafter\def\csname v@sim.0.fm3nw0\endcsname{5.0\%}
\expandafter\def\csname v@sim.0.fm3nw0.3\endcsname{1.0\%}
\expandafter\def\csname v@sim.0.hlA\endcsname{6.5\%}
\expandafter\def\csname v@sim.0.hlA.3\endcsname{1.0\%}
\expandafter\def\csname v@sim.0.hlAC\endcsname{6.5\%}
\expandafter\def\csname v@sim.0.hlAC.3\endcsname{1.0\%}
\expandafter\def\csname v@sim.0.hlResA\endcsname{7.0\%}
\expandafter\def\csname v@sim.0.hlResA.3\endcsname{0.0\%}
\expandafter\def\csname v@sim.0.reps\endcsname{200}
\expandafter\def\csname v@sim.P.fm2\endcsname{18\%}
\expandafter\def\csname v@sim.P.fm3\endcsname{100\%}
\expandafter\def\csname v@sim.P.fm3.t\endcsname{4.8}
\expandafter\def\csname v@sim.P.hlResA\endcsname{76\%}
\expandafter\def\csname v@sim.P.reps\endcsname{50}
\expandafter\def\csname v@sim.V.fm2\endcsname{100\%}
\expandafter\def\csname v@sim.V.fm2.t\endcsname{7.2}
\expandafter\def\csname v@sim.V.fm3\endcsname{4\%}
\expandafter\def\csname v@sim.V.hlAC\endcsname{100\%}
\expandafter\def\csname v@sim.V.hlResA\endcsname{8\%}
\expandafter\def\csname v@sim.V.reps\endcsname{50}
\expandafter\def\csname v@sim.assets\endcsname{300}
\expandafter\def\csname v@sim.cfg8\endcsname{3e8eb87e}
\expandafter\def\csname v@sim.checks.failed\endcsname{мир P, hl\_ResA: 76.0\% (при правиле ≥ 80\%); мир 0, fm3\_ResLnA: 9.5\% (95\% ДИ 5.8\%–14.4\% не содержит 5\%); мир 0, fm4\_ResLnA: 9.5\% (95\% ДИ 5.8\%–14.4\% не содержит 5\%)}
\expandafter\def\csname v@sim.checks.n\endcsname{17}
\expandafter\def\csname v@sim.checks.pass\endcsname{14}
\expandafter\def\csname v@sim.git\endcsname{aa8f4aa (dirty)}
\expandafter\def\csname v@sim.months\endcsname{90}
\expandafter\def\csname v@sim.run\endcsname{sim\_20261016\_191215}
\expandafter\def\csname v@sim.wall\endcsname{143}
\expandafter\def\csname v@t5.imla\endcsname{0.234 (t 0.50)}
\expandafter\def\csname v@t5.imla3\endcsname{0.155 (t 0.35)}
\expandafter\def\csname v@t5.raw\endcsname{1.54 (t 0.98)}
\expandafter\def\csname v@t7.jan\endcsname{0.699 (t 0.39)}
\expandafter\def\csname v@t7.jan.n\endcsname{7}
\expandafter\def\csname v@t7.rest\endcsname{0.766 (t 1.62)}
\expandafter\def\csname v@t7.rest.n\endcsname{80}
\expandafter\def\csname v@w0.crit_max\endcsname{2.22}
\expandafter\def\csname v@w0.crit_min\endcsname{2.00}
\expandafter\def\csname v@w0.exp_base\endcsname{5.7}
\expandafter\def\csname v@w0.fpr3_max\endcsname{1.5\%}
\expandafter\def\csname v@w0.fpr_max\endcsname{9.5\%}
\expandafter\def\csname v@w0.fpr_min\endcsname{6.5\%}
\expandafter\def\csname v@x21.fm2\endcsname{0.351 (t 0.97)}
\expandafter\def\csname v@x21.fm3\endcsname{−4.93 (t −2.54)}
\expandafter\def\csname v@x21.hlAC\endcsname{0.432 (t 0.42)}
\def\basedir{../results/20261016T135240Z_3e8eb87e}
\def\teamnames{Team 12}
\def\repourl{https://github.com/bigpurota/paper15-lou-shu}

\newcommand\V[1]{\ifcsname v@#1\endcsname\csname v@#1\endcsname\else\errmessage{undefined value: #1}\fi}

\definecolor{ink}{HTML}{1F2937}
\definecolor{accent}{HTML}{1D4E89}
\definecolor{muted}{HTML}{6B7280}
\definecolor{warn}{HTML}{9A3412}
\setbeamercolor{frametitle}{fg=accent}
\setbeamercolor{normal text}{fg=ink}
\setbeamercolor{structure}{fg=accent}
\setbeamerfont{frametitle}{size=\large,series=\bfseries}
\setbeamertemplate{navigation symbols}{}
\setbeamertemplate{itemize items}{\textbullet}
\setbeamertemplate{footline}{\hfill{\color{muted}\scriptsize Team 12 · Lou \& Shu (2017) on Binance \quad \insertframenumber\ /\ \inserttotalframenumber}\hspace{1.5em}\vspace{0.6em}}
\newcommand\msg[1]{{\color{accent}\textbf{#1}}\par\medskip}
\newcommand\src[1]{\par\vfill{\color{muted}\scriptsize Источник: #1}}
\renewcommand\arraystretch{1.15}

\begin{document}

% ------------------------------------------------------------------------------------------ 1
\begin{frame}[plain]
\vspace{1.5em}
{\color{muted}Team 12 · Paper 15 · RS3, DSBA-3, ВШЭ · защита 17.10.2026}\par\bigskip
{\Large\bfseries\color{accent} Price Impact or Trading Volume:\\ Why Is the Amihud (2002) Measure Priced?}\par\medskip
X. Lou, T. Shu, \emph{Review of Financial Studies} 30(12), 2017\par\bigskip
Репликация на Binance spot (пары к USDT, \V{n.first}…\V{n.last}) и апгрейды: Amihud \& Noh (2021), λ по 5-минутным барам, EDGE, сентимент, MOEX, США после 2012\par\vfill
\textbf{Команда:} \teamnames\par
\textbf{Репозиторий:} \ifx\repourl\empty{\color{warn}ссылка будет добавлена}\else\url{\repourl}\fi\par
{\color{muted}\scriptsize Все числа на слайдах — макросы из \texttt{results/} (src/build\_docs.py), базовый прогон \V{base.run}}
\end{frame}

% ------------------------------------------------------------------------------------------ 2
\begin{frame}{Problem: мера Amihud прайсится — но за что?}
\msg{Неликвидные по Amihud акции дают премию. Принято считать, что это плата за price impact. Lou \& Shu спрашивают, так ли это.}
\[ A_{i,t} = \frac{1}{D}\sum_{d} \frac{|r_{i,d}|}{\mathrm{DVol}_{i,d}} \]
\begin{itemize}
\item Премия у статьи: квинтильный H-L по $A$ — \V{pvu.HLArawCoef.paper} \%/мес (NYSE/AMEX 1964–2012, \V{d.paperMonths} месяцев).
\item Если премия за price impact, $A$ годится и чтобы \emph{мерить} ликвидность, и чтобы \emph{объяснять} доходности. Статья показывает, что это два разных вопроса.
\end{itemize}
\bigskip
{\small\color{muted}Место в литературе: мера (Amihud 2002) → проверка мер по HF-бенчмаркам (Goyenko–Holden–Trzcinka 2009) → за что прайсится (Lou \& Shu 2017) → критика декомпозиции (Amihud \& Noh 2021) → спред по OHLC (EDGE, Ardia–Guidotti–Kroencke 2024).}
\end{frame}

% ------------------------------------------------------------------------------------------ 3
\begin{frame}{Наблюдение: в поперечном сечении $A$ — это почти только объём}
\msg{Разброс объёма огромен, разброс $|r|$ мал, поэтому $A \approx A_C$.}
\begin{center}
\begin{tabular}{@{}lrr@{}}
\toprule
 & \textbf{статья (NYSE/AMEX)} & \textbf{мы (Binance)} \\
\midrule
corr($A$, $A_C$) & \V{pvu.corrAAC.paper} & \V{pvu.corrAAC} \\
Q75/Q25 для $A_C$ & \V{pvu.QratioAC.paper} & \V{pvu.QratioAC} \\
Q75/Q25 для $|Ret|$ & \V{pvu.QratioAbs.paper} & \V{pvu.QratioAbs} \\
corr($A$, $|Ret|$) & \V{pvu.corrAabs.paper} & \V{pvu.corrAabs} \\
\bottomrule
\end{tabular}
\end{center}
\begin{itemize}
\item Разброс ликвидности у нас на порядок меньше (фильтр объёма и сам листинг Binance — гипотеза), но вывод тот же: $A$ почти совпадает с $A_C$.
\item Отличие: у монет объём растёт с волатильностью, поэтому corr($A$, $|Ret|$) отрицательна.
\end{itemize}
\src{\V{base.dir}/paper\_vs\_us.csv (T1)}
\end{frame}

% ------------------------------------------------------------------------------------------ 4
\begin{frame}{Метод: «константная» мера и остаток}
\msg{Если прайсится $A_C$, а часть $A$ сверх $A_C$ — нет, то премия Amihud — это премия за объём.}
\[ A_C = \frac{1}{D}\sum_d \frac{1}{\mathrm{DVol}_d}, \qquad \ln A = a_t + b_t \ln A_C + \mathrm{Res.}\ln A \ \ \text{(в каждом месяце)} \]
\begin{columns}[T]
\begin{column}{0.48\textwidth}
\textbf{Гипотеза «объём»}\par прайсится $A_C$, остаток — нет
\end{column}
\begin{column}{0.48\textwidth}
\textbf{Гипотеза «price impact»}\par прайсится остаток Res.$A$
\end{column}
\end{columns}
\bigskip
\textbf{Сначала проверили, что конвейер различает гипотезы} (симуляции, \V{sim.assets} монет × \V{sim.months} месяцев):
\begin{itemize}
\item мир V (премия за объём): $A_C$ значим в \V{sim.V.fm2} повторов, Res.lnA — в \V{sim.V.fm3};
\item мир P (премия за price impact): Res.lnA значим в FM в \V{sim.P.fm3}, в квинтильной сортировке — только в \V{sim.P.hlResA}.
\end{itemize}
\src{results/simulations/sim\_world\_summary.csv}
\end{frame}

% ------------------------------------------------------------------------------------------ 5
\begin{frame}{Метод: тесты и данные}
\msg{Квинтильные сортировки и Fama–MacBeth с сигналом из $t-2$ и Newey–West(6).}
\centerline{мера за месяц $t-2$ \quad → \quad пропуск $t-1$ \quad → \quad доходность месяца $t$}
\[ r_{i,t} - \hat\beta_i\,\mathrm{CMKT}_t = \gamma_0 + \gamma_1 X_{i,t-2} + \text{controls} + \varepsilon_{i,t} \]
\begin{columns}[T]
\begin{column}{0.48\textwidth}
\textbf{Статья:} NYSE/AMEX 1964–2012, \V{d.paperMonths} месяцев, \V{d.paperAvgObs} акций в месяц; FF3-скорректированная доходность.
\end{column}
\begin{column}{0.48\textwidth}
\textbf{Мы:} Binance spot USDT, все пары с делистингом, \V{n.months} месяцев, в среднем \V{n.avgcoins} монет (\V{n.coins} всего); FM — \V{n.fmmonths} месяцев по \V{n.fmobs} монет. CMKT вместо FF3 (нет размера и B/M).
\end{column}
\end{columns}
\bigskip
\begin{itemize}
\item Аудит данных: \V{audit.redenomination} редоминаций скорректированы как сплиты, \V{audit.ticker_reuse} повторных тикера разведены (даты совпадают с остановками торгов).
\item Спецификация зафиксирована в git \V{git.prereg} до первого прогона на реальных данных; config первого прогона \V{git.prereg.cfgmatch} с предрегистрированным.
\end{itemize}
\end{frame}

% ------------------------------------------------------------------------------------------ 6
\begin{frame}{Результаты статьи (Table 3, Table 4)}
\msg{У статьи спред по $A$ и по $A_C$ положительный и значимый, по Res.$A$ — нет; в FM Res.lnA значимо отрицателен.}
\begin{center}
\begin{tabular}{@{}lcc@{}}
\toprule
\textbf{сигнал} & \textbf{квинтили Q1…Q5, \%/мес} & \textbf{H-L} \\
\midrule
$A$ & \V{q.A.paper} & \V{pvu.HLArawCoef.paper} \\
$A_C$ & \V{q.A_C.paper} & \V{pvu.HLACrawCoef.paper} \\
Res.$A$ & \V{q.ResA.paper} & \V{pvu.HLResArawCoef.paper} \\
\bottomrule
\end{tabular}
\end{center}
\begin{itemize}
\item FM: $\ln A_C$ \V{pvu.FMlnAC.paper}; Res.lnA \V{pvu.FMResLnA.paper}.
\item Логика статьи: часть $A$ сверх объёма прайсится со знаком минус, значит, премия не за price impact.
\end{itemize}
\src{src/paper\_numbers.yaml (сверено с PDF)}
\end{frame}

% ------------------------------------------------------------------------------------------ 7
\begin{frame}{Наша репликация рядом со статьёй}
\msg{Форма совпадает, но наши спреды шумные и незначимые.}
\begin{columns}[c]
\begin{column}{0.64\textwidth}
\includegraphics[width=\linewidth]{\basedir/fig_quintiles.png}
\end{column}
\begin{column}{0.34\textwidth}
\small
H-L $A_C$: \V{pvu.HLACrawCoef}\\ (статья \V{pvu.HLACrawCoef.paper})\par\smallskip
H-L $A$: \V{pvu.HLArawCoef}\par\smallskip
H-L Res.$A$: \V{pvu.HLResArawCoef}\par\medskip
s.d. месячной доходности монеты \V{d.retSd}\%.\par\smallskip
Без марта 2021 (альтсезон) спред $A_C$ \V{d.exMeanAC} \%/мес (t \V{d.exTAC}).
\end{column}
\end{columns}
\src{\V{base.dir}/fig\_quintiles.png: статья (Table 3, Panel A) и мы на одних осях}
\end{frame}

% ------------------------------------------------------------------------------------------ 8
\begin{frame}{Статья против нас}
\msg{Знаки совпадают у \V{sc.samesign} из \V{sc.ls_n} коэффициентов. Значимость повторяется только у Res.lnA (|t| ≥ 1.96 в \V{sc.res_sig} из \V{sc.res_n}, HLZ — ни разу); премия за объём незначима.}
{\scriptsize% generated by src/build_docs.py: do not edit
\begin{tabular}{@{}p{0.38\linewidth}p{0.21\linewidth}p{0.21\linewidth}p{0.11\linewidth}@{}}
\toprule
\textbf{metric} & \textbf{paper (NYSE/AMEX 1964–2012)} & \textbf{us (Binance)} & \textbf{match} \\
\midrule
corr(A, A\_C) & 0.899 & 0.956 & yes \\
Q75/Q25 of A\_C & 116.60 & 9.40 & no \\
corr(A, |Ret|) & 0.489 & −0.133 & no \\
H-L on A, raw (\%/mo) & 0.56 (t 2.36) & 1.54 (t 0.98) & sign only \\
H-L on A\_C, raw (\%/mo) & 0.61 (t 2.95) & 1.70 (t 1.12) & sign only \\
H-L on A\_C, alpha (\%/mo) & 0.44 (t 3.20) & 0.781 (t 0.62) & sign only \\
H-L on Res.A, raw (\%/mo) & −0.17 (t −1.05) & −0.510 (t −0.49) & yes \\
FM ln A, model (1) & 0.119 (t 3.01) & 0.580 (t 1.24) & sign only \\
FM ln A\_C, model (2) & 0.183 (t 4.79) & 0.760 (t 1.53) & sign only \\
FM Res.lnA, model (3) & −0.383 (t −5.07) & −5.09 (t −2.46) & yes \\
FM Res.lnA, model (4) with ln A\_C & −0.248 (t −3.14) & −8.10 (t −2.37) & yes \\
alpha of H-L(A) on IMLA\_C (\%/mo) & −0.04 (t −0.77) & 0.234 (t 0.50) & yes \\
\bottomrule
\end{tabular}
}
\par\smallskip
{\small Почему: (1) мощность — \V{n.months} месяцев против \V{d.paperMonths}, MDE80 спреда $A_C$ \V{mde.hlAC} \%/мес; (2) другой рынок — объём растёт с волатильностью; (3) нет размера и B/M, CMKT вместо FF3.}
\src{\V{base.dir}/paper\_vs\_us.md (все \V{sc.ls_n} строк с объяснениями — backup)}
\end{frame}

% ------------------------------------------------------------------------------------------ 9
\begin{frame}{HF-бенчмарк: λ по 5-минутным барам с точным знаком}
\msg{λ хорошо мерит $A$ и прайсится со знаком минус. С моделью (1) статьи это расходится, с моделью (2) — нет; премии $A_C$ в этом окне нет, Table 9 не воспроизведена.}
{\scriptsize% generated by src/build_docs.py: do not edit
\begin{tabular}{@{}p{0.33\linewidth}p{0.3\linewidth}p{0.31\linewidth}@{}}
\toprule
\textbf{} & \textbf{paper (T9, T10: λ sample 1983–2012, 358 mo.)} & \textbf{us (Binance 2023-10…2026-09, 5-min bars)} \\
\midrule
corr(λ, A) & 0.737 & 0.615 \\
corr(λ, A\_C) & 0.747 & 0.547 \\
corr(λ, Res.A) & 0.149 & 0.294 \\
FM ln λ, model (1) & −0.034 (t −0.67) & −2.59 (t −4.03) \\
FM ln λ with ln A\_C, model (2) & −0.155 (t −2.51) & −4.30 (t −3.57) \\
FM ln A\_C with ln λ, model (2) & 0.18 (t 3.80) & 1.06 (t 1.97) \\
T10: cost part, model (3) & 0.015 (t 0.35) & −1.42 (t −3.92) \\
T10: non-cost part, model (3) & 0.191 (t 3.10) & 0.957 (t 1.88) \\
\bottomrule
\end{tabular}
}
\par\smallskip
{\small
\begin{itemize}
\item \V{hf.months} месяцев, \V{hf.coinmonths} монето-месяцев; знак из taker-buy объёма 5-минутных klines. Бар-λ сверен с потрейдовым λ по aggTrades на \V{hf.val.n} монето-месяцах: corr ln λ \V{hf.val.corr}.
\item С контролем долларового объёма ln λ \V{hf.vol.b1}; ln $A_C$ без λ в этом окне \V{hf.b0}.
\end{itemize}}
\src{results/upgrades/hf\_lambda/ (T9A\_corr, T9B\_fm, T10\_fm). Статья: T9 — 1983–2012, T10 — 1964–2012}
\end{frame}

% ------------------------------------------------------------------------------------------ 10
\begin{frame}{Mispricing: сентимент, неликвидность рынка, январь}
\msg{Сентимент премию не объясняет; неликвидность рынка даёт знак статьи, но слабеет с контролем рынка.}
{\scriptsize% generated by src/build_docs.py: do not edit
\begin{tabular}{@{}p{0.33\linewidth}p{0.29\linewidth}p{0.32\linewidth}@{}}
\toprule
\textbf{test} & \textbf{paper (AT\_C; BW / PS)} & \textbf{us (A\_C; Fear \& Greed / AR(2) of A)} \\
\midrule
Sentiment → L−S (Table 13) & 0.24 (t 1.81) & 2.57 (t 1.84) \\
Sentiment → short leg & −0.25 (t −2.67) & −1.37 (t −1.70) \\
Sentiment → FM slope of ln A\_C & — & 0.586 (t 1.99) \\
Market illiquidity → L−S (Table 12) & −0.048 (t −2.38) & −4.30 (t −3.58) \\
same, + CMKT(t−1), post hoc & — & −3.16 (t −2.60) \\
January slope (Table 11 (1); paper: ln AT\_C) & −0.367 (t −2.27) & January 0.699 (t 0.39), n 7 \\
Other months & 0.23 (t 6.28) & 0.766 (t 1.62), n 80 \\
\bottomrule
\end{tabular}
}
\par\smallskip
{\small
\begin{itemize}
\item Сравнимы знак и t, не масштаб: у статьи сортировки по turnover и индексы Baker–Wurgler / Pástor–Stambaugh.
\item Сентимент: 1.96 проходят \V{se.var_196} из \V{se.var_n} предзаданных вариантов. Январей у нас \V{t7.jan.n}: январь в крипте — только плацебо.
\end{itemize}}
\src{results/upgrades/sentiment/paper\_vs\_us.csv, panel\_b.csv; \V{base.dir}/T7\_january.csv}
\end{frame}

% ------------------------------------------------------------------------------------------ 11
\begin{frame}{Кто прав: Lou \& Shu или Amihud \& Noh?}
\msg{Предпосылка A\&N на крипте выполняется, но их член DIF не прайсится; данные спор не разрешают.}
\begin{columns}[c]
\begin{column}{0.55\textwidth}
\includegraphics[width=\linewidth]{\basedir/fig_an_dif.png}
\end{column}
\begin{column}{0.43\textwidth}
\small
$\ln A = \ln|Ret| + \ln A_C + \mathrm{DIF}$\par\smallskip
COV < 0 в \V{an.covneg} монето-месяцев\par\smallskip
DIF при $\ln A_C$, $\ln|Ret|$: \V{an.V1m} (мес.), \V{an.V1y} (12 мес.); у A\&N \V{an.p.col3}\par\smallskip
fDIF (12 мес.): \V{an.V4y} — единственный результат в сторону A\&N, ниже HLZ\par\smallskip
Res.lnA: \V{an.V3m} → \V{an.V5m} при $\ln|Ret|$ + IVOL. Но это почти тождество: R² Res.lnA на $\ln|Ret|$ + DIF = \V{an.r2_res_dif}. «Остаток — это волатильность» — гипотеза.
\end{column}
\end{columns}
\src{results/upgrades/amihud\_noh/; вердикт правила: \V{an.sens.ret_adj} при минимуме монет \V{an.sens.levels}}
\end{frame}

% ------------------------------------------------------------------------------------------ 12
\begin{frame}{Критика}
\msg{У статьи слабое звено — интерпретация остатка; у нас — мощность и один месяц.}
\begin{columns}[c]
\begin{column}{0.46\textwidth}
\includegraphics[width=\linewidth]{\basedir/fig_hl_cumulative.png}
\end{column}
\begin{column}{0.52\textwidth}
\footnotesize
\begin{enumerate}
\item Декомпозиция теряет ковариационный член (A\&N); остаток ≈ $|Ret|$ (IVOL-пазл, Ang et al. 2006).
\item «Объём, а не ликвидность» — ложная дихотомия: ликвидные активы больше торгуются (Amihud–Mendelson 1986).
\item Шумный бенчмарк — λ статьи (5 минут, Lee–Ready): ошибка измерения занижает его коэффициент.
\item EW-портфели, только США до 2012: после 2012 премия Illiquidity исчезла (OpenAP: \V{op.ill.pre} → \V{op.ill.post}).
\end{enumerate}
У нас: без 2021 года спред $A_C$ \V{x21.hlAC}; на MOEX спред поглощается размером (SMB \V{mx.smb}), а FM $\ln A_C$ \V{mx.fm2} повторяет статью.
\end{column}
\end{columns}
\end{frame}

% ------------------------------------------------------------------------------------------ 13
\begin{frame}{Переобучение: могли ли результаты быть переобучены?}
\small
\textbf{Авторы:} подпериоды, NASDAQ вне выборки, годовые меры, сырая доходность.\par\medskip
\textbf{Мы:}
\begin{itemize}
\item Спецификация закоммичена до первого прогона (git \V{git.prereg}); реестр всех t: в базовом семействе \V{lg.base.n} тестов, |t| > 1.96 у \V{lg.base.196} (ожидалось бы \V{lg.base.exp}–\V{w0.exp_base} при нуле), |t| > 3 у \V{lg.base.3}, после Бонферрони/BH — \V{lg.base.bonf}/\V{lg.base.bh}.
\item Размер теста в мире 0: NW(6) даёт \V{w0.fpr_min}–\V{w0.fpr_max} ложных |t| > 1.96; эмпирический 5\%-порог \V{w0.crit_min}–\V{w0.crit_max}. Симуляции: \V{sim.checks.pass} из \V{sim.checks.n} критериев; не пройдены: {\color{warn}\V{sim.checks.failed}}.
\item Бутстрап по блокам месяцев: ДИ без нуля устойчиво только у \V{rb.robust} из \V{rb.n} статистик (обе — Res.lnA).
\item Плацебо: перемешивание мер внутри месяца — средний t по \V{plc.shuffle_reps} перестановкам не больше \V{plc.shuffle_maxabs_meant} по модулю, доля |t| > 1.96 не выше \V{plc.shuffle_share196_max}; утечка (лаг 0) даёт |t| до \V{plc.leak_maxabs_t}.
\item Вне выборки: другой рынок (MOEX), США после 2012 (OpenAP), без 2021 года. Post hoc варианты помечены.
\end{itemize}
\src{results/upgrades/robustness/ (ledger\_summary, bootstrap\_ci, world0\_calibration); \V{base.dir}/placebo.csv}
\end{frame}

% ------------------------------------------------------------------------------------------ 14
\begin{frame}{Next steps: что бы мы сделали за 3 месяца}
\msg{Главное — мощность и измерение спреда.}
\begin{enumerate}
\item Капитализация монет (supply из CoinGecko/CMC) → контроль размера, $AT$/$AT_C$ и VW-портфели.
\item Внутридневной EDGE на 5-минутных барах и спреды из снапшотов стакана.
\item λ по midquote и на всей истории, а не на \V{hf.months} месяцах.
\item Тест A\&N на 12-месячных мерах с заранее заданным минимумом монет; остаток на $\ln$ VOLUME, как в Table 6.
\item McLean–Pontiff: распад премии после публикации внутри выборки США.
\end{enumerate}
\bigskip
{\small MDE80 для FM $\ln A_C$: \V{mde.fm2.slope} на единицу $\ln A_C$ против наклона статьи \V{mde.fm2.paper_slope} (в \V{mde.fm2.slope_ratio} раза больше). Нужно больше месяцев или меньше шума, а не больше спецификаций.}
\end{frame}

% ========================================================================================== backup
\begin{frame}[plain]
\centering\vfill{\Large\color{accent}\bfseries Backup}\vfill
\end{frame}

\begin{frame}{B1. Определения}
\small
\begin{itemize}
\item $A = \mathrm{mean}_d(|r_d|/\mathrm{DVol}_d)$, $A_C = \mathrm{mean}_d(1/\mathrm{DVol}_d)$; Res.lnA — остаток $\ln A$ на $\ln A_C$ в каждом месяце.
\item Amihud \& Noh: $\ln A = \ln|Ret| + \ln A_C + \mathrm{DIF}$, DIF отражает $\mathrm{cov}(|r|, 1/\mathrm{DVol})$.
\item λ (Hasbrouck): 5-минутная доходность на $\sum \mathrm{sign}\cdot\sqrt{\mathrm{DVol}}$, у нас SVol $= \sqrt{B} - \sqrt{S}$ по taker-объёму.
\item IMLA$_C$ — терцильный H-L по $A_C$; CMKT — EW-рынок вселенной; CMOM — терцильный H-L по Ret[−12,−2].
\end{itemize}
\src{src/measures.py, src/hf\_lambda.py, src/factors.py}
\end{frame}

\begin{frame}{B2. Симуляции: миры V / P / 0}
\msg{В мире V прайсится $A_C$, в мире P — Res.lnA; в мире 0 NW(6) слегка переотвергает.}
{\scriptsize% generated by src/build_docs.py: do not edit
\begin{tabular}{@{}p{0.05\linewidth}p{0.16\linewidth}p{0.33\linewidth}p{0.12\linewidth}p{0.09\linewidth}p{0.09\linewidth}@{}}
\toprule
\textbf{мир} & \textbf{статистика} & \textbf{ожидание} & \textbf{доля значимых} & \textbf{средний t} & \textbf{повторов} \\
\midrule
V & fm2\_lnA\_C & прайсится: доля t>1.96 высокая & 100.0\% & 7.24 & 50 \\
V & hl\_A\_C & прайсится & 100.0\% & 6.61 & 50 \\
V & fm3\_ResLnA & не прайсится: доля |t|>1.96 около 5\% & 4.0\% & 0.06 & 50 \\
V & hl\_ResA & не прайсится & 8.0\% & −0.05 & 50 \\
P & fm3\_ResLnA & прайсится & 100.0\% & 4.78 & 50 \\
P & hl\_ResA & прайсится (сортировка слабее FM) & 76.0\% & 2.83 & 50 \\
P & fm2\_lnA\_C & в идеале не прайсится & 18.0\% & 0.01 & 50 \\
0 & fm2\_lnA\_C & ложные значимости около 5\% & 8.0\% & −0.03 & 200 \\
0 & fm3\_ResLnA & ложные значимости около 5\% & 9.5\% & −0.14 & 200 \\
0 & fm3\_ResLnA\_nw0 & то же без NW-лагов & 5.0\% & −0.12 & 200 \\
0 & hl\_A\_C & ложные значимости около 5\% & 6.5\% & −0.06 & 200 \\
\bottomrule
\end{tabular}
}
\par\medskip
{\small Пройдено \V{sim.checks.pass} из \V{sim.checks.n} критериев. Не пройдены: \V{sim.checks.failed}.}
\src{results/simulations/sim\_world\_summary.csv, sim\_checks.csv}
\end{frame}

\begin{frame}{B3. t-статистики статьи и наши}
\centering\includegraphics[height=0.8\textheight]{\basedir/fig_tstats.png}
\src{\V{base.dir}/fig\_tstats.png (линии 1.96 и HLZ 3)}
\end{frame}

\begin{frame}{B4. Мощность и устойчивость}
\small
\begin{itemize}
\item MDE80: H-L $A_C$ \V{mde.hlAC} \%/мес против премии статьи \V{mde.hlAC.paper}. FM $\ln A_C$: \V{mde.fm2.slope} на единицу $\ln A_C$ против \V{mde.fm2.paper_slope}; на 1 s.d. (каждое в своём s.d.) \V{mde.fm2} (s.d. \V{mde.fm2.sd}) против \V{mde.fm2.paper} (s.d. статьи \V{mde.fm2.paper_sd}).
\item Без 2021 года: H-L $A_C$ \V{x21.hlAC}, FM $\ln A_C$ \V{x21.fm2}, Res.lnA (3) \V{x21.fm3}.
\item \V{rob.same} из \V{rob.n} вариантов (винзоризация, сырая доходность, веса, подпериоды) сохраняют вывод. Подпериоды: $\ln A_C$ \V{pvu.SubAlnAC} и \V{pvu.SubBlnAC}.
\item Бутстрап: FM Res.lnA (3) \V{rb.ci.fm3_ResLnA}, (4) \V{rb.ci.fm4_ResLnA}; FM $\ln A_C$ (4) на грани (\V{rb.seedshare.fm4_lnA_C} сидов без нуля).
\end{itemize}
\src{results/upgrades/robustness/; \V{base.dir}/robustness\_vs\_base.csv}
\end{frame}

\begin{frame}{B5. Остальные апгрейды}
\small
\begin{itemize}
\item \textbf{EDGE и издержки:} H-L $A_C$ брутто \V{ed.gross}, нетто \V{ed.net}; дневной EDGE спред не мерит (вес гэпа \V{ed.w2}); безубыточная издержка \V{ed.breakeven} на сторону.
\item \textbf{MOEX TQBR:} FM $\ln A$ \V{mx.fm1}, $\ln A_C$ \V{mx.fm2}, Res.lnA (4) \V{mx.fm4res}; спред $A_C$ \V{mx.hlAC}, альфа \V{mx.hlACalpha}; \V{mx.topyear} год даёт \V{mx.topyear_share} суммарного спреда.
\item \textbf{США (OpenAP):} Illiquidity \V{op.ill.pre} → \V{op.ill.post}, DolVol \V{op.dv.pre} → \V{op.dv.post}; t разницы по лагам NW \V{op.ill.decay_lags}: \V{op.ill.decay_ts}.
\item \textbf{Реестр тестов:} базовые тесты и апгрейды вместе — \V{lg.all.n} тестов, |t| > 3 у \V{lg.all.3}, после Бонферрони \V{lg.all.bonf}.
\end{itemize}
\src{results/upgrades/*/SUMMARY.md}
\end{frame}

\begin{frame}{B6. Воспроизводимость}
\small
\begin{itemize}
\item \texttt{make all}: данные → тесты → симуляции → базовый прогон → отчёт → апгрейды → реестр → документы; \texttt{make slides} собирает эту колоду.
\item Каждый прогон пишет \texttt{run\_manifest.json}: run id, git SHA и флаг dirty, sha256 конфига, кода и данных, seed, время по этапам.
\item Отступления от статьи: DEVIATIONS.md (D1–D17, U1–U8, журнал J1–J2).
\end{itemize}
\src{README.md}
\end{frame}

\end{document}
